\subsection{PRODUCT OF UNIFORM SPACES}

DEFINITION 4. If \((\mathbf{X}_t)_{i\in I}\) is a family of uniform spaces, the product uniform space of this family is the product set

\[
\mathbf {X} = \prod_ {i \in I} \mathbf {X} _ {i}
\]

endowed with the coarsest uniformity for which the projections \(pr_{t}: X \rightarrow X_{t}\) are uniformly continuous. This uniformity is called the product of the uniformities of the \(X_{t}\) , and the uniform spaces \(X_{t}\) are called the factors of X.

The topology induced by the product uniformity on X is same as the product of the topologies of the \(X_{i}\) (no. 3, Corollary to Proposition 4).

PROPOSITION 7. Let \(f = (f_{t})\) be a mapping of a uniform space \(Y\) into a product uniform space \(X = \prod_{i \in I} X_{i}\) . Then \(f\) is uniformly continuous if and only if each \(f_{t}\) is uniformly continuous.

Since \(f_{i} = \mathrm{pr}_{i} \circ f\) , this is a particular case of Proposition 4 of no. 3.

COROLLARY. Let \((\mathbf{X}_t)_{i \in I}, (\mathbf{Y}_t)_{i \in I}\) be two families of uniform spaces indexed by the same set I. For each \(i \in I\) , let \(f_t\) be a mapping of \(\mathbf{X}_t\) into \(\mathbf{Y}_t\) . If each of the \(f_t\) is uniformly continuous, then so is the product mapping

\[
f: (x _ {i}) \rightarrow (f _ {i} (x _ {i})).
\]

Conversely, if the \(X_{i}\) are non-empty and f is uniformly continuous, then each \(f_{i}\) is uniformly continuous.

\(f\) can be written as \(x \to (f_{i}(\mathrm{pr}_{i}x))\) , and the first part of the corollary therefore follows from Proposition 7. The second part is proved by considering a point \(a = (a_{i})\) of \(\prod_{i \in I} X_{i}\) and repeating the argument of Chapter I, § 4, no. 1, Corollary I to Proposition 1, with the phrase ``continuous at \(a\) (resp. \(a_{x}\) )'' replaced by ``uniformly continuous''.

The general criterion of transitivity of initial uniformities (no. 3, Proposition 5) shows that, as for the product of topological spaces (Chapter I, § 4, no. 1), the product of uniform spaces is associative and that the following is true:

PROPOSITION 8. Let X be a set, let \((\mathbf{Y}_t)_{t \in I}\) be a family of uniform spaces, and for each \(t \in I\) let \(f_t\) be a mapping of X into \(\mathbf{Y}_t\) . Let \(f\) be the mapping \(x \to (f_t(x))\) of X into \(\mathbf{Y} = \prod_{t \in I} \mathbf{Y}_t\) , and let \(\mathfrak{U}\) be the coarsest uniformity on X for which the \(f_t\) are uniformly continuous. Then \(\mathfrak{U}\) is the inverse image under \(f\) of the uniformity induced on \(f(\mathbf{X})\) by the product uniformity on Y.

COROLLARY. For each \(i \in I\) , let \(A_{i}\) be a subspace of \(Y_{i}\) . Then the uniformity induced on \(A = \prod_{i \in I} A_{i}\) by the product uniformity on \(\prod_{i \in I} Y_{i}\) is the same as the product of the uniformities of the subspaces \(A_{i}\) .

In addition, we see immediately that if \(X_{1}, X_{2}\) are two uniform spaces and \(a_{1}\) is any point of \(X_{1}\) , the mapping \(x_{2} \to (a_{1}, x_{2})\) is an isomorphism of \(X_{2}\) onto the subspace \(\{a_{1}\} \times X_{2}\) of \(X_{1} \times X_{2}\) ; hence:

PROPOSITION 9. Let \(f\) be a uniformly continuous mapping of a product uniform space \(\mathbf{X}_1 \times \mathbf{X}_2\) into a uniform space \(\mathbf{Y}\) ; then every partial mapping

\[
x _ {2} \rightarrow f (x _ {1}, x _ {2})
\]

of \(X_{2}\) into Y is uniformly continuous.

In other words, a uniformly continuous function of two arguments is uniformly continuous with respect to each of them separately.

* The example given in Chapter I, § 4, no. 2, Remark 2, shows that the converse of this proposition is false. *

